\section{Conclusions}
\label{sec:conclusions}

We validated an analytic covariance model for the DESI DR2 LRG1 and QSO power-spectrum multipoles against 859 approximate and 25 $N$-body mocks. The main results:
\begin{enumerate}
\item The Gaussian term is accurate for the window it is given. For QSO it agrees with the mocks to $\simeq1$\% in $\langle\chi^2\rangle/n$, per-element variances and parameter errors.
\item For LRG1 the local window $m^2$ underestimates the Gaussian variance by $\simeq5$\%, because the veto mask varies on scales below the correlation length. A window pair-averaged with the $\xi$-kernel recovers it with no free parameter. Object-to-object weight variations must be averaged locally, not taken per random, for the same reason.
\item The remaining LRG1 excess is concentrated in two amplitude modes of the monopole and quadrupole, the signature of super-sample covariance. Super-sample covariance with the local-average effect, together with the window-convolved discreteness terms and nothing fitted, brings $\langle\chi^2\rangle/n$ to $\val{trace/LRG1/NGC/rec/0.3}$ (NGC) and $\val{trace/LRG1/SGC/rec/0.3}$ (SGC) at $k<0.3\,\hMpc$, and parameter-error ratios to within 10\% except for the quadrupole amplitude.
\item The perturbative connected trispectrum, at tree level and in the response approximation, is rejected by both mock suites at the predicted amplitude.
\item Fibre assignment does not change the covariance beyond the change of density and weights included in the window.
\item The residual is a 7--16\% excess of the hexadecapole variance conditioned on the lower multipoles, and a 12--16\% excess of the quadrupole-amplitude variance. Its origin is open.
\end{enumerate}

On methodology:
\begin{itemize}
\item Distributions of $\chi^2$ and per-element variance ratios are weak tests of a covariance matrix. Most of the information is in a few directions, which must be found without overfitting: in-sample eigenvalues of the whitened sample covariance overstate the excess by a factor of two, while cross-validated variances are unbiased.
\item Varying the bin width separates Gaussian-like errors from missing non-Gaussian terms, and an eigenvalue that does not change with the bin width identifies a low-rank term.
\item A validation against mocks is only as complete as the mocks. The approximate mocks lack the long radial modes that dominate super-sample covariance in a thin shell, and this should be checked before the agreement of such a term is taken as validation.
\end{itemize}
